% =====================================================================================
% SELF-CONTAINED: the bibliography is embedded below via filecontents, so this single
% .tex file is all you need to paste into Overleaf. On compile LaTeX writes it out as
% iclr2027_conference.bib and BibTeX picks it up. If you would rather keep the .bib as a
% separate file in the project, delete the filecontents block -- nothing else changes.
% =====================================================================================
\begin{filecontents*}[overwrite]{iclr2027_conference.bib}
% =====================================================================================
% References for "How Gradient Descent Builds a Noise-Skipping Induction Head".
%
% PROVENANCE, so this can be audited:
%  [V] author list verified against arXiv (source: Insights/_make_related_papers_pdf.py,
%      "Papers verified against arXiv API on 2026-07-04"), or a widely-known reference.
%  [P] PARTIAL -- lead author(s) verified, remainder recorded as "and others". Expand from
%      the paper before submission.
%  [C] classic / foundational, cited from common knowledge; check the exact venue string.
%
% The ICLR 2027 AI policy makes the authors responsible for every citation. Every [P]
% entry below MUST be completed and every entry spot-checked before submission.
% =====================================================================================

% ------------------------------------------------------------------ induction heads [V/C]
@article{elhage2021,
  title   = {A Mathematical Framework for Transformer Circuits},
  author  = {Elhage, Nelson and Nanda, Neel and Olsson, Catherine and Henighan, Tom and
             Joseph, Nicholas and Mann, Ben and Askell, Amanda and Bai, Yuntao and
             Chen, Anna and Conerly, Tom and DasSarma, Nova and Drain, Dawn and
             Ganguli, Deep and Hatfield-Dodds, Zac and Hernandez, Danny and
             Jones, Andy and Kernion, Jackson and Lovitt, Liane and Ndousse, Kamal and
             Amodei, Dario and Brown, Tom and Clark, Jack and Kaplan, Jared and
             McCandlish, Sam and Olah, Chris},
  journal = {Transformer Circuits Thread},
  year    = {2021}
}

@article{olsson2022,
  title   = {In-context Learning and Induction Heads},
  author  = {Olsson, Catherine and Elhage, Nelson and Nanda, Neel and Joseph, Nicholas and
             DasSarma, Nova and Henighan, Tom and Mann, Ben and Askell, Amanda and
             Bai, Yuntao and Chen, Anna and Conerly, Tom and Drain, Dawn and
             Ganguli, Deep and Hatfield-Dodds, Zac and Hernandez, Danny and
             Johnston, Scott and Jones, Andy and Kernion, Jackson and Lovitt, Liane and
             Ndousse, Kamal and Amodei, Dario and Brown, Tom and Clark, Jack and
             Kaplan, Jared and McCandlish, Sam and Olah, Chris},
  journal = {Transformer Circuits Thread},
  year    = {2022}
}

@inproceedings{bietti2023,
  title     = {Birth of a Transformer: A Memory Viewpoint},
  author    = {Bietti, Alberto and Cabannes, Vivien and Bouchacourt, Diane and
               J{\'e}gou, Herv{\'e} and Bottou, L{\'e}on},
  booktitle = {Advances in Neural Information Processing Systems},
  year      = {2023},
  note      = {arXiv:2306.00802}
}

% -------------------------------------------------- in-context learning of Markov structure
@inproceedings{edelman2024,
  title     = {The Evolution of Statistical Induction Heads: In-Context Learning Markov Chains},
  author    = {Edelman, Ezra and Tsilivis, Nikolaos and Edelman, Benjamin L. and
               Malach, Eran and Goel, Surbhi},
  booktitle = {Advances in Neural Information Processing Systems},
  year      = {2024},
  note      = {arXiv:2402.11004}
}

@inproceedings{ekbote2025,
  title     = {What One Cannot, Two Can: Two-Layer Transformers Provably Represent
               Induction Heads on Any-Order Markov Chains},
  author    = {Ekbote, Chanakya and Bondaschi, Marco and Rajaraman, Nived and Lee, Jason D.
               and Gastpar, Michael and Makkuva, Ashok Vardhan and Liang, Paul Pu},
  booktitle = {Advances in Neural Information Processing Systems},
  year      = {2025},
  note      = {arXiv:2508.07208}
}

@article{shi2025,                                                            % [P]
  title   = {Towards Understanding Transformers in Learning Random Walks},
  author  = {Shi, Yiwei and Cao, Yuan},
  journal = {arXiv preprint arXiv:2511.23239},
  year    = {2025},
  note    = {[P] verify given names}
}

@article{dai2026,                                                            % [P]
  title   = {Pre-trained Large Language Models Learn Hidden Markov Models In-context},
  author  = {Dai, Yijia and others},
  journal = {arXiv preprint arXiv:2506.07298},
  year    = {2026},
  note    = {[P] expand author list}
}

@article{hao2025,                                                            % [P]
  title   = {Transformers as Multi-task Learners: Decoupling Features in Hidden Markov Models},
  author  = {Hao, Yifan and others},
  journal = {arXiv preprint arXiv:2506.01919},
  year    = {2025},
  note    = {[P] expand author list}
}

% ------------------------------------------------------------------- training dynamics [V/P]
@inproceedings{nichani2024,
  title     = {How Transformers Learn Causal Structure with Gradient Descent},
  author    = {Nichani, Eshaan and Damian, Alex and Lee, Jason D.},
  booktitle = {International Conference on Machine Learning},
  year      = {2024},
  note      = {arXiv:2402.14735}
}

@inproceedings{unveiling2024,
  title     = {Unveiling Induction Heads: Provable Training Dynamics and Feature Learning
               in Transformers},
  author    = {Chen, Siyu and Sheen, Heejune and Wang, Tianhao and Yang, Zhuoran},
  booktitle = {Advances in Neural Information Processing Systems},
  year      = {2024},
  note      = {arXiv:2409.10559}
}

@article{varre2025,                                                          % [P]
  title   = {Learning In-context $n$-grams with Transformers: Sub-$n$-grams Are
             Near-stationary Points},
  author  = {Varre, Aditya and others},
  journal = {arXiv preprint arXiv:2508.12837},
  year    = {2025},
  note    = {[P] expand author list}
}

@article{musat2026,                                                          % [P]
  title   = {On the Emergence of Induction Heads for In-Context Learning},
  author  = {Musat, Andrei and others},
  journal = {arXiv preprint arXiv:2511.01033},
  year    = {2026},
  note    = {[P] expand author list}
}

@article{wang2025,                                                       % [P]
  title   = {How Transformers Get Rich: Approximation and Dynamics Analysis},
  author  = {Wang, Mingze and others},
  journal = {arXiv preprint arXiv:2410.11474},
  year    = {2025},
  note    = {[P] expand author list}
}

% ------------------------------------------------ data distribution, shortcuts, emergence [P]
@inproceedings{kawata2025,
  title     = {From Shortcut to Induction Head: How Data Diversity Shapes Algorithm
               Selection in Transformers},
  author    = {Kawata, Yuma and others},
  booktitle = {Advances in Neural Information Processing Systems},
  year      = {2025},
  note      = {arXiv:2512.18634. [P] expand author list}
}

@inproceedings{zucchet2025,
  title     = {The Emergence of Sparse Attention: Impact of Data Distribution and
               Benefits of Repetition},
  author    = {Zucchet, Nicolas and others},
  booktitle = {Advances in Neural Information Processing Systems},
  year      = {2025},
  note      = {arXiv:2505.17863. [P] expand author list}
}

@article{lavie2026,                                                          % [P]
  title   = {Phase Transitions in Attention: A Bayesian Theory of Copy Head Emergence},
  author  = {Lavie, Ori and others},
  journal = {arXiv preprint arXiv:2606.12058},
  year    = {2026},
  note    = {[P] expand author list}
}

@article{yuksel2026,                                                         % [P]
  title   = {Incremental Learning of Sparse Attention Patterns in Transformers},
  author  = {Yuksel, Kamer Ali and others},
  journal = {arXiv preprint arXiv:2602.19143},
  year    = {2026},
  note    = {[P] expand author list}
}

% ------------------------------------------------------ selecting causal structure in context
@inproceedings{selective2025,
  title     = {Selective Induction Heads: How Transformers Select Causal Structures in Context},
  author    = {Anonymous},
  booktitle = {International Conference on Learning Representations},
  year      = {2025},
  note      = {arXiv:2509.08184. [P] REQUIRED: fill in the author list from the paper}
}

% ------------------------------------------------------------ Bayesian views of ICL [P]
@article{agarwal2026,
  title   = {The Bayesian Geometry of Transformer Attention},
  author  = {Agarwal, Rohan and others},
  journal = {arXiv preprint arXiv:2512.22471},
  year    = {2026},
  note    = {companion: Gradient Dynamics of Attention, arXiv:2512.22473. [P] expand}
}

@article{wakayama2025,                                                       % [P]
  title   = {In-Context Learning Is Provably Bayesian Inference: A Generalization Theory
             for Meta-Learning},
  author  = {Wakayama, Tomoya and Imaizumi, Masaaki},
  journal = {arXiv preprint arXiv:2510.10981},
  year    = {2025},
  note    = {[P] verify given names}
}

@article{zhang2026,                                                          % [P]
  title   = {Bayesian Optimality of In-Context Learning with Selective State Spaces},
  author  = {Zhang, Yifan and others},
  journal = {arXiv preprint arXiv:2602.17744},
  year    = {2026},
  note    = {[P] expand author list}
}

@article{dangelo2026,                                                        % [P]
  title   = {Transformers Learn Latent Mixture Models In-Context via Mirror Descent},
  author  = {D'Angelo, Francesco and Flammarion, Nicolas},
  journal = {arXiv preprint arXiv:2604.10848},
  year    = {2026},
  note    = {[P] verify given names}
}

% ---------------------------------------------------- denoising, FFN, architecture [P]
@inproceedings{smart2025,
  title     = {In-context Denoising with One-layer Transformers: Connections between
               Attention and Associative Memory Retrieval},
  author    = {Smart, Matthew and others},
  booktitle = {International Conference on Machine Learning},
  year      = {2025},
  note      = {arXiv:2502.05164. [P] expand author list}
}

@article{sun2025,                                                            % [P]
  title   = {On the Role of Transformer Feed-Forward Layers in Nonlinear In-Context Learning},
  author  = {Sun, Haoyuan and others},
  journal = {arXiv preprint arXiv:2501.18187},
  year    = {2025},
  note    = {[P] expand author list}
}

% ------------------------------------------------- symmetry, gauge, parameterisation [P]
@article{silverstein2026,
  title   = {Symmetry Breaking in Transformers for Efficient and Interpretable Training},
  author  = {Silverstein, Louis and others},
  journal = {arXiv preprint arXiv:2601.22257},
  year    = {2026},
  note    = {[P] expand author list}
}

@inproceedings{kobayashi2024,                                                % [P]
  title     = {Weight Decay Induces Low-rank Attention Layers},
  author    = {Kobayashi, Seijin and others},
  booktitle = {Advances in Neural Information Processing Systems},
  year      = {2024},
  note      = {arXiv:2410.23819. [P] expand author list}
}

% ------------------------------------------------------------------- foundational [C]
@inproceedings{vaswani2017,
  title     = {Attention Is All You Need},
  author    = {Vaswani, Ashish and Shazeer, Noam and Parmar, Niki and Uszkoreit, Jakob and
               Jones, Llion and Gomez, Aidan N. and Kaiser, Lukasz and Polosukhin, Illia},
  booktitle = {Advances in Neural Information Processing Systems},
  year      = {2017}
}

@inproceedings{brown2020,
  title     = {Language Models Are Few-Shot Learners},
  author    = {Brown, Tom B. and Mann, Benjamin and Ryder, Nick and Subbiah, Melanie and
               Kaplan, Jared and Dhariwal, Prafulla and Neelakantan, Arvind and others},
  booktitle = {Advances in Neural Information Processing Systems},
  year      = {2020},
  note      = {[P] expand author list if cited prominently}
}

@inproceedings{garg2022,
  title     = {What Can Transformers Learn In-Context? A Case Study of Simple Function Classes},
  author    = {Garg, Shivam and Tsipras, Dimitris and Liang, Percy and Valiant, Gregory},
  booktitle = {Advances in Neural Information Processing Systems},
  year      = {2022}
}

@inproceedings{xie2022,
  title     = {An Explanation of In-context Learning as Implicit Bayesian Inference},
  author    = {Xie, Sang Michael and Raghunathan, Aditi and Liang, Percy and Ma, Tengyu},
  booktitle = {International Conference on Learning Representations},
  year      = {2022}
}

@inproceedings{vonoswald2023,
  title     = {Transformers Learn In-Context by Gradient Descent},
  author    = {von Oswald, Johannes and Niklasson, Eyvind and Randazzo, Ettore and
               Sacramento, Jo{\~a}o and Mordvintsev, Alexander and Zhmoginov, Andrey and
               Vladymyrov, Max},
  booktitle = {International Conference on Machine Learning},
  year      = {2023}
}

@inproceedings{akyurek2023,
  title     = {What Learning Algorithm Is In-Context Learning? Investigations with Linear Models},
  author    = {Aky{\"u}rek, Ekin and Schuurmans, Dale and Andreas, Jacob and Ma, Tengyu and
               Zhou, Denny},
  booktitle = {International Conference on Learning Representations},
  year      = {2023}
}

\end{filecontents*}

\documentclass{article} % For LaTeX2e
\usepackage{iclr2027_conference,times}

% Optional math commands from https://github.com/goodfeli/dlbook_notation.
%%%%% NEW MATH DEFINITIONS %%%%%

\usepackage{amsmath,amsfonts,bm}

% Mark sections of captions for referring to divisions of figures
\newcommand{\figleft}{{\em (Left)}}
\newcommand{\figcenter}{{\em (Center)}}
\newcommand{\figright}{{\em (Right)}}
\newcommand{\figtop}{{\em (Top)}}
\newcommand{\figbottom}{{\em (Bottom)}}
\newcommand{\captiona}{{\em (a)}}
\newcommand{\captionb}{{\em (b)}}
\newcommand{\captionc}{{\em (c)}}
\newcommand{\captiond}{{\em (d)}}

% Highlight a newly defined term
\newcommand{\newterm}[1]{{\bf #1}}


% Figure reference, lower-case.
\def\figref#1{figure~\ref{#1}}
% Figure reference, capital. For start of sentence
\def\Figref#1{Figure~\ref{#1}}
\def\twofigref#1#2{figures \ref{#1} and \ref{#2}}
\def\quadfigref#1#2#3#4{figures \ref{#1}, \ref{#2}, \ref{#3} and \ref{#4}}
% Section reference, lower-case.
\def\secref#1{section~\ref{#1}}
% Section reference, capital.
\def\Secref#1{Section~\ref{#1}}
% Reference to two sections.
\def\twosecrefs#1#2{sections \ref{#1} and \ref{#2}}
% Reference to three sections.
\def\secrefs#1#2#3{sections \ref{#1}, \ref{#2} and \ref{#3}}
% Reference to an equation, lower-case.
\def\eqref#1{equation~\ref{#1}}
% Reference to an equation, upper case
\def\Eqref#1{Equation~\ref{#1}}
% A raw reference to an equation---avoid using if possible
\def\plaineqref#1{\ref{#1}}
% Reference to a chapter, lower-case.
\def\chapref#1{chapter~\ref{#1}}
% Reference to an equation, upper case.
\def\Chapref#1{Chapter~\ref{#1}}
% Reference to a range of chapters
\def\rangechapref#1#2{chapters\ref{#1}--\ref{#2}}
% Reference to an algorithm, lower-case.
\def\algref#1{algorithm~\ref{#1}}
% Reference to an algorithm, upper case.
\def\Algref#1{Algorithm~\ref{#1}}
\def\twoalgref#1#2{algorithms \ref{#1} and \ref{#2}}
\def\Twoalgref#1#2{Algorithms \ref{#1} and \ref{#2}}
% Reference to a part, lower case
\def\partref#1{part~\ref{#1}}
% Reference to a part, upper case
\def\Partref#1{Part~\ref{#1}}
\def\twopartref#1#2{parts \ref{#1} and \ref{#2}}

\def\ceil#1{\lceil #1 \rceil}
\def\floor#1{\lfloor #1 \rfloor}
\def\1{\bm{1}}
\newcommand{\train}{\mathcal{D}}
\newcommand{\valid}{\mathcal{D_{\mathrm{valid}}}}
\newcommand{\test}{\mathcal{D_{\mathrm{test}}}}

\def\eps{{\epsilon}}


% Random variables
\def\reta{{\textnormal{$\eta$}}}
\def\ra{{\textnormal{a}}}
\def\rb{{\textnormal{b}}}
\def\rc{{\textnormal{c}}}
\def\rd{{\textnormal{d}}}
\def\re{{\textnormal{e}}}
\def\rf{{\textnormal{f}}}
\def\rg{{\textnormal{g}}}
\def\rh{{\textnormal{h}}}
\def\ri{{\textnormal{i}}}
\def\rj{{\textnormal{j}}}
\def\rk{{\textnormal{k}}}
\def\rl{{\textnormal{l}}}
% rm is already a command, just don't name any random variables m
\def\rn{{\textnormal{n}}}
\def\ro{{\textnormal{o}}}
\def\rp{{\textnormal{p}}}
\def\rq{{\textnormal{q}}}
\def\rr{{\textnormal{r}}}
\def\rs{{\textnormal{s}}}
\def\rt{{\textnormal{t}}}
\def\ru{{\textnormal{u}}}
\def\rv{{\textnormal{v}}}
\def\rw{{\textnormal{w}}}
\def\rx{{\textnormal{x}}}
\def\ry{{\textnormal{y}}}
\def\rz{{\textnormal{z}}}

% Random vectors
\def\rvepsilon{{\mathbf{\epsilon}}}
\def\rvtheta{{\mathbf{\theta}}}
\def\rva{{\mathbf{a}}}
\def\rvb{{\mathbf{b}}}
\def\rvc{{\mathbf{c}}}
\def\rvd{{\mathbf{d}}}
\def\rve{{\mathbf{e}}}
\def\rvf{{\mathbf{f}}}
\def\rvg{{\mathbf{g}}}
\def\rvh{{\mathbf{h}}}
\def\rvu{{\mathbf{i}}}
\def\rvj{{\mathbf{j}}}
\def\rvk{{\mathbf{k}}}
\def\rvl{{\mathbf{l}}}
\def\rvm{{\mathbf{m}}}
\def\rvn{{\mathbf{n}}}
\def\rvo{{\mathbf{o}}}
\def\rvp{{\mathbf{p}}}
\def\rvq{{\mathbf{q}}}
\def\rvr{{\mathbf{r}}}
\def\rvs{{\mathbf{s}}}
\def\rvt{{\mathbf{t}}}
\def\rvu{{\mathbf{u}}}
\def\rvv{{\mathbf{v}}}
\def\rvw{{\mathbf{w}}}
\def\rvx{{\mathbf{x}}}
\def\rvy{{\mathbf{y}}}
\def\rvz{{\mathbf{z}}}

% Elements of random vectors
\def\erva{{\textnormal{a}}}
\def\ervb{{\textnormal{b}}}
\def\ervc{{\textnormal{c}}}
\def\ervd{{\textnormal{d}}}
\def\erve{{\textnormal{e}}}
\def\ervf{{\textnormal{f}}}
\def\ervg{{\textnormal{g}}}
\def\ervh{{\textnormal{h}}}
\def\ervi{{\textnormal{i}}}
\def\ervj{{\textnormal{j}}}
\def\ervk{{\textnormal{k}}}
\def\ervl{{\textnormal{l}}}
\def\ervm{{\textnormal{m}}}
\def\ervn{{\textnormal{n}}}
\def\ervo{{\textnormal{o}}}
\def\ervp{{\textnormal{p}}}
\def\ervq{{\textnormal{q}}}
\def\ervr{{\textnormal{r}}}
\def\ervs{{\textnormal{s}}}
\def\ervt{{\textnormal{t}}}
\def\ervu{{\textnormal{u}}}
\def\ervv{{\textnormal{v}}}
\def\ervw{{\textnormal{w}}}
\def\ervx{{\textnormal{x}}}
\def\ervy{{\textnormal{y}}}
\def\ervz{{\textnormal{z}}}

% Random matrices
\def\rmA{{\mathbf{A}}}
\def\rmB{{\mathbf{B}}}
\def\rmC{{\mathbf{C}}}
\def\rmD{{\mathbf{D}}}
\def\rmE{{\mathbf{E}}}
\def\rmF{{\mathbf{F}}}
\def\rmG{{\mathbf{G}}}
\def\rmH{{\mathbf{H}}}
\def\rmI{{\mathbf{I}}}
\def\rmJ{{\mathbf{J}}}
\def\rmK{{\mathbf{K}}}
\def\rmL{{\mathbf{L}}}
\def\rmM{{\mathbf{M}}}
\def\rmN{{\mathbf{N}}}
\def\rmO{{\mathbf{O}}}
\def\rmP{{\mathbf{P}}}
\def\rmQ{{\mathbf{Q}}}
\def\rmR{{\mathbf{R}}}
\def\rmS{{\mathbf{S}}}
\def\rmT{{\mathbf{T}}}
\def\rmU{{\mathbf{U}}}
\def\rmV{{\mathbf{V}}}
\def\rmW{{\mathbf{W}}}
\def\rmX{{\mathbf{X}}}
\def\rmY{{\mathbf{Y}}}
\def\rmZ{{\mathbf{Z}}}

% Elements of random matrices
\def\ermA{{\textnormal{A}}}
\def\ermB{{\textnormal{B}}}
\def\ermC{{\textnormal{C}}}
\def\ermD{{\textnormal{D}}}
\def\ermE{{\textnormal{E}}}
\def\ermF{{\textnormal{F}}}
\def\ermG{{\textnormal{G}}}
\def\ermH{{\textnormal{H}}}
\def\ermI{{\textnormal{I}}}
\def\ermJ{{\textnormal{J}}}
\def\ermK{{\textnormal{K}}}
\def\ermL{{\textnormal{L}}}
\def\ermM{{\textnormal{M}}}
\def\ermN{{\textnormal{N}}}
\def\ermO{{\textnormal{O}}}
\def\ermP{{\textnormal{P}}}
\def\ermQ{{\textnormal{Q}}}
\def\ermR{{\textnormal{R}}}
\def\ermS{{\textnormal{S}}}
\def\ermT{{\textnormal{T}}}
\def\ermU{{\textnormal{U}}}
\def\ermV{{\textnormal{V}}}
\def\ermW{{\textnormal{W}}}
\def\ermX{{\textnormal{X}}}
\def\ermY{{\textnormal{Y}}}
\def\ermZ{{\textnormal{Z}}}

% Vectors
\def\vzero{{\bm{0}}}
\def\vone{{\bm{1}}}
\def\vmu{{\bm{\mu}}}
\def\vtheta{{\bm{\theta}}}
\def\va{{\bm{a}}}
\def\vb{{\bm{b}}}
\def\vc{{\bm{c}}}
\def\vd{{\bm{d}}}
\def\ve{{\bm{e}}}
\def\vf{{\bm{f}}}
\def\vg{{\bm{g}}}
\def\vh{{\bm{h}}}
\def\vi{{\bm{i}}}
\def\vj{{\bm{j}}}
\def\vk{{\bm{k}}}
\def\vl{{\bm{l}}}
\def\vm{{\bm{m}}}
\def\vn{{\bm{n}}}
\def\vo{{\bm{o}}}
\def\vp{{\bm{p}}}
\def\vq{{\bm{q}}}
\def\vr{{\bm{r}}}
\def\vs{{\bm{s}}}
\def\vt{{\bm{t}}}
\def\vu{{\bm{u}}}
\def\vv{{\bm{v}}}
\def\vw{{\bm{w}}}
\def\vx{{\bm{x}}}
\def\vy{{\bm{y}}}
\def\vz{{\bm{z}}}

% Elements of vectors
\def\evalpha{{\alpha}}
\def\evbeta{{\beta}}
\def\evepsilon{{\epsilon}}
\def\evlambda{{\lambda}}
\def\evomega{{\omega}}
\def\evmu{{\mu}}
\def\evpsi{{\psi}}
\def\evsigma{{\sigma}}
\def\evtheta{{\theta}}
\def\eva{{a}}
\def\evb{{b}}
\def\evc{{c}}
\def\evd{{d}}
\def\eve{{e}}
\def\evf{{f}}
\def\evg{{g}}
\def\evh{{h}}
\def\evi{{i}}
\def\evj{{j}}
\def\evk{{k}}
\def\evl{{l}}
\def\evm{{m}}
\def\evn{{n}}
\def\evo{{o}}
\def\evp{{p}}
\def\evq{{q}}
\def\evr{{r}}
\def\evs{{s}}
\def\evt{{t}}
\def\evu{{u}}
\def\evv{{v}}
\def\evw{{w}}
\def\evx{{x}}
\def\evy{{y}}
\def\evz{{z}}

% Matrix
\def\mA{{\bm{A}}}
\def\mB{{\bm{B}}}
\def\mC{{\bm{C}}}
\def\mD{{\bm{D}}}
\def\mE{{\bm{E}}}
\def\mF{{\bm{F}}}
\def\mG{{\bm{G}}}
\def\mH{{\bm{H}}}
\def\mI{{\bm{I}}}
\def\mJ{{\bm{J}}}
\def\mK{{\bm{K}}}
\def\mL{{\bm{L}}}
\def\mM{{\bm{M}}}
\def\mN{{\bm{N}}}
\def\mO{{\bm{O}}}
\def\mP{{\bm{P}}}
\def\mQ{{\bm{Q}}}
\def\mR{{\bm{R}}}
\def\mS{{\bm{S}}}
\def\mT{{\bm{T}}}
\def\mU{{\bm{U}}}
\def\mV{{\bm{V}}}
\def\mW{{\bm{W}}}
\def\mX{{\bm{X}}}
\def\mY{{\bm{Y}}}
\def\mZ{{\bm{Z}}}
\def\mBeta{{\bm{\beta}}}
\def\mPhi{{\bm{\Phi}}}
\def\mLambda{{\bm{\Lambda}}}
\def\mSigma{{\bm{\Sigma}}}

% Tensor
\DeclareMathAlphabet{\mathsfit}{\encodingdefault}{\sfdefault}{m}{sl}
\SetMathAlphabet{\mathsfit}{bold}{\encodingdefault}{\sfdefault}{bx}{n}
\newcommand{\tens}[1]{\bm{\mathsfit{#1}}}
\def\tA{{\tens{A}}}
\def\tB{{\tens{B}}}
\def\tC{{\tens{C}}}
\def\tD{{\tens{D}}}
\def\tE{{\tens{E}}}
\def\tF{{\tens{F}}}
\def\tG{{\tens{G}}}
\def\tH{{\tens{H}}}
\def\tI{{\tens{I}}}
\def\tJ{{\tens{J}}}
\def\tK{{\tens{K}}}
\def\tL{{\tens{L}}}
\def\tM{{\tens{M}}}
\def\tN{{\tens{N}}}
\def\tO{{\tens{O}}}
\def\tP{{\tens{P}}}
\def\tQ{{\tens{Q}}}
\def\tR{{\tens{R}}}
\def\tS{{\tens{S}}}
\def\tT{{\tens{T}}}
\def\tU{{\tens{U}}}
\def\tV{{\tens{V}}}
\def\tW{{\tens{W}}}
\def\tX{{\tens{X}}}
\def\tY{{\tens{Y}}}
\def\tZ{{\tens{Z}}}


% Graph
\def\gA{{\mathcal{A}}}
\def\gB{{\mathcal{B}}}
\def\gC{{\mathcal{C}}}
\def\gD{{\mathcal{D}}}
\def\gE{{\mathcal{E}}}
\def\gF{{\mathcal{F}}}
\def\gG{{\mathcal{G}}}
\def\gH{{\mathcal{H}}}
\def\gI{{\mathcal{I}}}
\def\gJ{{\mathcal{J}}}
\def\gK{{\mathcal{K}}}
\def\gL{{\mathcal{L}}}
\def\gM{{\mathcal{M}}}
\def\gN{{\mathcal{N}}}
\def\gO{{\mathcal{O}}}
\def\gP{{\mathcal{P}}}
\def\gQ{{\mathcal{Q}}}
\def\gR{{\mathcal{R}}}
\def\gS{{\mathcal{S}}}
\def\gT{{\mathcal{T}}}
\def\gU{{\mathcal{U}}}
\def\gV{{\mathcal{V}}}
\def\gW{{\mathcal{W}}}
\def\gX{{\mathcal{X}}}
\def\gY{{\mathcal{Y}}}
\def\gZ{{\mathcal{Z}}}

% Sets
\def\sA{{\mathbb{A}}}
\def\sB{{\mathbb{B}}}
\def\sC{{\mathbb{C}}}
\def\sD{{\mathbb{D}}}
% Don't use a set called E, because this would be the same as our symbol
% for expectation.
\def\sF{{\mathbb{F}}}
\def\sG{{\mathbb{G}}}
\def\sH{{\mathbb{H}}}
\def\sI{{\mathbb{I}}}
\def\sJ{{\mathbb{J}}}
\def\sK{{\mathbb{K}}}
\def\sL{{\mathbb{L}}}
\def\sM{{\mathbb{M}}}
\def\sN{{\mathbb{N}}}
\def\sO{{\mathbb{O}}}
\def\sP{{\mathbb{P}}}
\def\sQ{{\mathbb{Q}}}
\def\sR{{\mathbb{R}}}
\def\sS{{\mathbb{S}}}
\def\sT{{\mathbb{T}}}
\def\sU{{\mathbb{U}}}
\def\sV{{\mathbb{V}}}
\def\sW{{\mathbb{W}}}
\def\sX{{\mathbb{X}}}
\def\sY{{\mathbb{Y}}}
\def\sZ{{\mathbb{Z}}}

% Entries of a matrix
\def\emLambda{{\Lambda}}
\def\emA{{A}}
\def\emB{{B}}
\def\emC{{C}}
\def\emD{{D}}
\def\emE{{E}}
\def\emF{{F}}
\def\emG{{G}}
\def\emH{{H}}
\def\emI{{I}}
\def\emJ{{J}}
\def\emK{{K}}
\def\emL{{L}}
\def\emM{{M}}
\def\emN{{N}}
\def\emO{{O}}
\def\emP{{P}}
\def\emQ{{Q}}
\def\emR{{R}}
\def\emS{{S}}
\def\emT{{T}}
\def\emU{{U}}
\def\emV{{V}}
\def\emW{{W}}
\def\emX{{X}}
\def\emY{{Y}}
\def\emZ{{Z}}
\def\emSigma{{\Sigma}}

% entries of a tensor
% Same font as tensor, without \bm wrapper
\newcommand{\etens}[1]{\mathsfit{#1}}
\def\etLambda{{\etens{\Lambda}}}
\def\etA{{\etens{A}}}
\def\etB{{\etens{B}}}
\def\etC{{\etens{C}}}
\def\etD{{\etens{D}}}
\def\etE{{\etens{E}}}
\def\etF{{\etens{F}}}
\def\etG{{\etens{G}}}
\def\etH{{\etens{H}}}
\def\etI{{\etens{I}}}
\def\etJ{{\etens{J}}}
\def\etK{{\etens{K}}}
\def\etL{{\etens{L}}}
\def\etM{{\etens{M}}}
\def\etN{{\etens{N}}}
\def\etO{{\etens{O}}}
\def\etP{{\etens{P}}}
\def\etQ{{\etens{Q}}}
\def\etR{{\etens{R}}}
\def\etS{{\etens{S}}}
\def\etT{{\etens{T}}}
\def\etU{{\etens{U}}}
\def\etV{{\etens{V}}}
\def\etW{{\etens{W}}}
\def\etX{{\etens{X}}}
\def\etY{{\etens{Y}}}
\def\etZ{{\etens{Z}}}

% The true underlying data generating distribution
\newcommand{\pdata}{p_{\rm{data}}}
% The empirical distribution defined by the training set
\newcommand{\ptrain}{\hat{p}_{\rm{data}}}
\newcommand{\Ptrain}{\hat{P}_{\rm{data}}}
% The model distribution
\newcommand{\pmodel}{p_{\rm{model}}}
\newcommand{\Pmodel}{P_{\rm{model}}}
\newcommand{\ptildemodel}{\tilde{p}_{\rm{model}}}
% Stochastic autoencoder distributions
\newcommand{\pencode}{p_{\rm{encoder}}}
\newcommand{\pdecode}{p_{\rm{decoder}}}
\newcommand{\precons}{p_{\rm{reconstruct}}}

\newcommand{\laplace}{\mathrm{Laplace}} % Laplace distribution

\newcommand{\E}{\mathbb{E}}
\newcommand{\Ls}{\mathcal{L}}
\newcommand{\R}{\mathbb{R}}
\newcommand{\emp}{\tilde{p}}
\newcommand{\lr}{\alpha}
\newcommand{\reg}{\lambda}
\newcommand{\rect}{\mathrm{rectifier}}
\newcommand{\softmax}{\mathrm{softmax}}
\newcommand{\sigmoid}{\sigma}
\newcommand{\softplus}{\zeta}
\newcommand{\KL}{D_{\mathrm{KL}}}
\newcommand{\Var}{\mathrm{Var}}
\newcommand{\standarderror}{\mathrm{SE}}
\newcommand{\Cov}{\mathrm{Cov}}
% Wolfram Mathworld says $L^2$ is for function spaces and $\ell^2$ is for vectors
% But then they seem to use $L^2$ for vectors throughout the site, and so does
% wikipedia.
\newcommand{\normlzero}{L^0}
\newcommand{\normlone}{L^1}
\newcommand{\normltwo}{L^2}
\newcommand{\normlp}{L^p}
\newcommand{\normmax}{L^\infty}

\newcommand{\parents}{Pa} % See usage in notation.tex. Chosen to match Daphne's book.

\DeclareMathOperator*{\argmax}{arg\,max}
\DeclareMathOperator*{\argmin}{arg\,min}

\DeclareMathOperator{\sign}{sign}
\DeclareMathOperator{\Tr}{Tr}
\let\ab\allowbreak


\usepackage{hyperref}
\usepackage{url}
\usepackage{graphicx}
\usepackage{booktabs}
\usepackage{amsmath,amssymb,amsthm}

\newtheorem{theorem}{Theorem}
\newtheorem{lemma}[theorem]{Lemma}
\newtheorem{proposition}[theorem]{Proposition}

% ---------------------------------------------------------------------------
% TODO before submitting:
%   * Keep \iclrfinalcopy COMMENTED OUT. A non-anonymous submission is
%     desk-rejected without review (ICLR 2027 author guidelines).
%   * Main text limit: 9 pages. References are UNLIMITED and do not count.
%     Appendices are unlimited. The AI Use Statement, Ethics Statement and
%     Reproducibility Statement do NOT count toward the page limit.
%   * Produce US Letter, not A4.
%   * Do not alter margins, font sizes or any style-file parameter.
% ---------------------------------------------------------------------------

% Title shortlist is kept in CLAUDE.md section 10.7.
% NOTE: the ICLR style sets the title as a JUSTIFIED full-width paragraph in \LARGE small caps
% (see \@maketitle in iclr2027_conference.sty). A title longer than roughly 40 characters wraps,
% and LaTeX then stretches the inter-word spaces of the first line to justify it. Control the
% break yourself with \\ -- the template's own title does exactly this. Do NOT edit the .sty.
\title{How Gradient Descent Builds a\\Noise-Skipping Induction Head}

% Authors are hidden while \iclrfinalcopy is commented out.
\author{Anonymous}

%\iclrfinalcopy % Uncomment for camera-ready version, but NOT for submission.
\begin{document}

\maketitle

\begin{abstract}
% ICLR requires the abstract to be a SINGLE paragraph.
Induction heads explain how transformers infer causal structure from context rather than from
their weights, but existing analyses place the informative token adjacent to the query, or at one
of a small enumerable set of candidate lags. We study in-context learning when the causal parent
is instead separated from its child by an unbounded run of uninformative material. Sequences are
drawn from a Markov chain whose transition table is resampled for every sequence, so nothing about
it can be stored in the weights, and a tunable fraction of positions emit a symbol from a disjoint
distractor sub-vocabulary that leaves the chain state unchanged. Prediction then requires an
additional primitive we call state recovery: the model must first skip backwards over a run of
distractors of unbounded length to locate the last state-bearing symbol, and only then can the
induction match be posed. Neither a previous-token head nor a selector over finitely many lags can
express this. We give a two-layer construction that solves the task, in which the first layer
performs state recovery and the second performs counting induction on the recovered state, and we
then analyse how gradient descent builds it. Starting from zero initialisation, where the model
computes the in-context unigram predictor and every gradient is available in closed form, we prove
that the first layer's gradient vanishes identically, so the two layers cannot form
simultaneously; that the second layer's gradient admits a closed form rewarding keys whose symbol
matches the target; and that, because the transition table is resampled, the channel corresponding
to weight-stored bigram statistics receives no signal in expectation. In-context learning is
therefore not merely preferred over in-weights learning in this setting but is the only direction
available to gradient descent at the start of training. The second layer then recruits the first,
which is driven towards state recovery, so the circuit is assembled outwards from the readout
rather than inwards from the input. The analysis further predicts that state recovery is learned
only when distractors are sufficiently common, since the positions that demand it are exactly
those the distractor process creates; we confirm this empirically, observing a sharp transition
beyond which the model abandons state recovery in favour of copying the previous symbol. As the
distractor rate vanishes the task, the construction and the dynamics all reduce to the noiseless
setting studied in prior work, which we recover as a boundary case.
\end{abstract}

\section{Introduction}
\label{sec:intro}

A language model shown a short passage can often continue it correctly even when the pattern it
must follow appears nowhere in its training data. The pattern is supplied by the passage itself,
and the model applies it on the spot. This ability --- in-context learning \citep{brown2020} --- has been
studied through the function classes a transformer can acquire from a prompt \citep{garg2022},
through its interpretation as implicit Bayesian inference \citep{xie2022}, and through the
algorithms attention can emulate internally \citep{vonoswald2023,akyurek2023}. It is now
understood to rest, at least in part, on a specific and identifiable circuit. An \emph{induction head}
\citep{elhage2021,olsson2022} scans the context for an earlier occurrence of the current symbol
and copies whatever followed it. Two attention layers suffice: the first records, at every
position, the identity of the symbol that preceded it; the second searches for positions whose
recorded predecessor matches the symbol at hand, and averages what it finds. The circuit is simple
enough to write down by hand, and simple enough that one can ask how gradient descent arrives at
it.

That question has been studied in settings built to isolate it. The usual construction draws each
sequence from a Markov chain whose transition table is resampled every time, so that nothing about
the table can be stored in the weights and all of the achievable performance must come from the
context \citep{edelman2024,nichani2024,unveiling2024}. In such settings the relevant evidence is
always close at hand: the symbol that determines the next one is the immediately preceding symbol,
or a parent at a known offset, or one of a handful of candidate lags among which the model
chooses \citep{selective2025}.

\paragraph{When the evidence is not adjacent.}
Sequences encountered in practice are rarely so orderly. Between two causally linked symbols there
may be an arbitrary amount of material that says nothing about the link: filler, formatting,
digression, interruption. A mechanism that can only compare against the previous position, or
against a fixed menu of positions, is structurally unable to see across such material. What it
does instead is the question this paper is about.

We therefore study the simplest task in which the difficulty is unavoidable. As before, each
sequence is generated by a Markov chain with its own freshly drawn transition table. In addition,
each position independently emits --- with probability $1-\alpha$ --- a symbol from a
\emph{distractor} sub-vocabulary disjoint from the chain's own alphabet, and such a symbol
\emph{does not advance the chain}. A symbol and its causal parent are consequently separated by a
run of distractors whose length is geometric with mean $(1-\alpha)^{-1}$, and which is not bounded
by any constant. Before the induction match can even be posed, the model must first carry out what
we call \emph{state recovery}: skip backwards across that run to find the last symbol that
actually moved the chain. The single parameter $\alpha$ tunes how severe this is, and at
$\alpha=1$ the distractors disappear and the task becomes exactly the one previously
studied.\footnote{The two axes are complementary rather than competing. \citet{selective2025} vary
the dependency distance while holding the transition structure fixed; we resample the structure
and leave the distance to be discovered. A selector over $k$ candidate lags cannot represent a
geometric lag distribution for any finite $k$.}

\paragraph{How the circuit gets built.}
We first confirm that the task is solvable in this architecture by exhibiting a two-layer
construction (Section~\ref{sec:construction}): the first layer vetoes distractor keys and applies
a recency preference among the survivors, so that it retrieves the most recent state-bearing
symbol at whatever distance it lies; the second performs the familiar induction match on that
recovered state. We then turn to the question the paper is named for --- how gradient descent
assembles this --- and analyse it where it is tractable, at initialisation.

Initialising both attention matrices to zero makes every attention distribution uniform, so the
model begins by predicting the running frequency of the symbols it has seen: the in-context
unigram predictor. Every gradient at that point is available in closed form, and three facts
follow (Theorem~\ref{thm:dynamics}). First, the gradient with respect to the first layer is
\emph{identically zero}: every path from its parameters to the loss passes through the second
layer's attention scores, and that derivative carries the second layer's weights as a
multiplicative factor. The two layers cannot form at the same time. Second, the second layer's
gradient has a closed form, and it raises the score of exactly those positions whose symbol equals
the target. Third --- and this is what ties the mechanism to the task --- when that update is
decomposed into interpretable blocks, the block that would implement a weight-stored bigram table
receives no structured signal at all, because averaging over the resampled table makes every
signal symbol equally likely to follow every other. In-weights learning is not merely worse here
than in-context learning; it is invisible to the gradient. The only direction available at the
first step is already an in-context one.

What happens next follows from the first fact. Once the second layer has moved, it \emph{recruits}
the first: the first layer's gradient becomes proportional to precisely those blocks of the second
layer that consume the recovered state, and it is driven to suppress distractors and sharpen onto
the most recent state-bearing position. The circuit is assembled outwards from the readout rather
than inwards from the input.

\paragraph{A prediction, and its confirmation.}
This account carries a consequence. The pressure to build state recovery originates only at
positions where the state is not directly visible --- those that follow a distractor --- and such
positions are a fraction $1-\alpha$ of the sequence. As distractors become rarer, that pressure
weakens, while a shortcut remains available: simply treat the previous symbol as the state, which
is correct a fraction $\alpha$ of the time. We find that the transition between the two regimes is
abrupt. Below a threshold in $\alpha$ the trained model almost never errs about which symbols may
follow the current state; above it, the model is wrong about half the time, which is close to what
ignoring the state entirely would give.\footnote{The direction of the effect is worth emphasising.
The regime in which the model fails is the one with \emph{less} noise, where the task is easier by
every aggregate measure: both the cross-entropy and its Bayes floor are lower. The failure is
confined to the one sub-task that distinguishes this setting from the noiseless one, and the
aggregate objective barely weights it.}

\paragraph{Contributions.}
\begin{itemize}
\item We formulate in-context learning of Markov structure \emph{under distraction}, in which the
causal parent lies at an unbounded distance, and give a two-layer construction solving it that
introduces state recovery as an explicit primitive (Section~\ref{sec:construction}).
\item We analyse gradient descent from zero initialisation and prove three properties of the first
step: the layers form sequentially, the second layer's gradient is available in closed form, and
the in-weights channel receives no signal in expectation (Section~\ref{sec:dynamics}; proofs in
Appendix~\ref{app:dynamics}).
\item We report a sharp phase transition in distractor density, anticipated by that analysis, in
which state recovery ceases to be learned once distractors become rare
(Section~\ref{sec:transition}).
\item We show that the task, the construction and the dynamics all reduce to the noiseless setting
as $\alpha\to1$, recovering prior results as a boundary case rather than competing with them
(Section~\ref{sec:reduction}).
\end{itemize}

\section{Setup}
\label{sec:setup}

\paragraph{Task.}
Fix a vocabulary $\mathcal{V}=\mathcal{S}\cup\mathcal{D}$ partitioned into $V_s$ \emph{signal}
symbols and $V_d$ \emph{distractors}. For each sequence we draw a transition table $\pi$
independently, assigning to every $s\in\mathcal{S}$ a successor set
$\mathrm{succ}_\pi(s)\subset\mathcal{S}$ of size two, and generate $v_{0:T-1}$ by
\begin{equation}
v_t=\begin{cases}
\text{uniform from }\mathrm{succ}_\pi(\mathrm{st}(t\!-\!1)), & \text{w.p. }\alpha
\quad\text{(state advances)},\\[2pt]
\text{uniform from }\mathcal{D}, & \text{w.p. }1-\alpha
\quad\text{(state unchanged)},
\end{cases}
\end{equation}
where $\mathrm{st}(t)$ denotes the chain state after position $t$. The target is
$y_t=v_{t+1}$ and the loss is averaged over the second half of the sequence,
$t\in\mathcal{T}=[T/2,T)$, so that every scored position has a context long enough to identify
the table. Because $\pi$ is resampled, no function of the weights alone can predict better than
the marginal; all of the achievable performance is in-context.

\paragraph{The state and why it is not the previous symbol.}
Distractors do not advance the chain, so the state at position $t$ is the most recent
\emph{signal} symbol at or before $t$. The number of distractors separating a symbol from its
causal parent is geometric with mean $(1-\alpha)^{-1}$ and is unbounded. A previous-token head,
or a selector over any finite set of lags, cannot express this dependency.

\paragraph{Architecture.}
The input embedding is a concatenation of two one-hots and is not learned,
$x_t=[\,e_{v_t}\,;\,e_t\,]\in\mathbb{R}^{d}$ with $d=V+T$. Two attention layers concatenate their
output onto the residual stream,\footnote{Concatenation rather than addition keeps the two layers'
contributions in disjoint coordinates, which is what makes the block decomposition of
Section~\ref{sec:dynamics} exact rather than approximate.} and both use a strictly causal mask.
With combined query--key matrices $W^{(1)}\in\mathbb{R}^{d\times d}$,
$W^{(2)}\in\mathbb{R}^{2d\times 2d}$,
\begin{equation}
h_t=\sum_{u<t} a^{(1)}_{t,u}\,e_{v_u},\qquad
x^{(2)}_t=[\,x_t\,;\,h_t\,],\qquad
f_t=\sum_{u<t} a^{(2)}_{t,u}\,e_{v_u},
\label{eq:arch}
\end{equation}
where $a^{(i)}_{t,\cdot}=\mathrm{softmax}\bigl(x^{(i)}_tW^{(i)}x^{(i)\top}_{\cdot}/\sqrt d+M\bigr)$
and the value circuits are fixed to the selector that passes the symbol one-hot. Appendix
\ref{app:gauge} shows that fixing the first layer's value circuit costs no generality --- it is a
choice of coordinates --- while fixing the second layer's is a genuine restriction, which we
state as such.

\paragraph{Readout.}
Since $f_t$ in \eqref{eq:arch} is a convex combination of one-hots it is already a distribution,
with $\sum_w f_t[w]=1$ identically. We therefore read out by renormalisation against a fixed
bias, with no unembedding, no logarithm and no output softmax:
\begin{equation}
p_t=\frac{f_t+b}{\sum_w\bigl(f_t[w]+b_w\bigr)},\qquad
b_w=\begin{cases} b_s, & w\in\mathcal{S},\\ b_d, & w\in\mathcal{D},\end{cases}
\qquad
b_d=\frac{(1-\alpha)(1+V_s b_s)}{V_d\,\alpha}.
\label{eq:readout}
\end{equation}
The choice of $b_d$ is forced: it is the unique value for which a model with a perfect distractor
filter assigns exactly $(1-\alpha)/V_d$ to each distractor. Appendix~\ref{app:readout} derives it
and shows that this readout is what an unembedding followed by a softmax computes when the
intervening nonlinearity is $\log(f+b_s)$.

\begin{table}[!ht]
\caption{Notation.}
\label{tab:notation}
\begin{center}
\begin{tabular}{ll}
\toprule
$\mathcal{S},\mathcal{D}$ & signal and distractor sub-vocabularies, sizes $V_s$, $V_d$ \\
$\alpha$ & probability that a position emits a signal symbol and advances the chain \\
$\pi$ & per-sequence transition table; $\mathrm{succ}_\pi(s)$ its successor set \\
$\mathrm{st}(t)$ & chain state after position $t$; the most recent signal symbol \\
$h_t, f_t$ & output of layer 1 and layer 2 at position $t$, both in the simplex \\
$m_t$ & running symbol frequency $\tfrac1t\sum_{u<t}e_{v_u}$ of the context \\
$P,Q,R,C$ & the four $d\times d$ blocks of $W^{(2)}$ (Section~\ref{sec:dynamics}) \\
\bottomrule
\end{tabular}
\end{center}
\end{table}

\section{The state-recovering induction head}
\label{sec:construction}

We first establish that the task is solvable by this architecture, and how. The construction
divides the work between the layers.

\begin{theorem}[Construction]
\label{thm:construction}
There exist weights $W^{(1)},W^{(2)}$ such that, up to softmax saturation, layer 1 places all of
its attention on the most recent signal position and layer 2 outputs the empirical successor
frequency of the current state:
\begin{equation}
h_t=e_{\mathrm{st}(t)},\qquad
f_t[w]=\frac{\#\{u<t:\ \mathrm{st}(u\!-\!1)=\mathrm{st}(t),\ v_u=w\}}
{\#\{u<t:\ \mathrm{st}(u\!-\!1)=\mathrm{st}(t)\}} .
\end{equation}
Composed with the readout \eqref{eq:readout}, the model realises add-$b_s$ smoothed
context-Bayes prediction.
\end{theorem}

The mechanism, proved in Appendix~\ref{app:construction}, uses three devices. First, layer 1
assigns score $-\infty$ to every key whose symbol is a distractor, so the attention is supported
on signal positions only. Second, among those it applies a recency ramp linear in $t-u$, so the
softmax concentrates on the most recent one; the ramp is what makes the unbounded lag expressible
with a fixed number of parameters. Third, layer 2 compares the query's own symbol against the key's
recovered state $h_u$, and averages the symbols it finds --- turning an attention average into a
count, since the attended set is exactly the set of past transitions out of the current state.

Two remarks. The first layer is \emph{not} a previous-token head: it is a previous-\emph{signal}-token
head, and the distinction is the whole content of the construction. And the query in layer 2
matches on its own symbol when that symbol is a signal, falling back to the recovered state when it
is a distractor; a single head suffices because those two branches ask the same question.

\section{A distractor-density phase transition}
\label{sec:transition}

Theorem~\ref{thm:construction} says the circuit exists. It does not say gradient descent finds it,
and the central empirical result of this paper is that whether it does depends sharply on
$\alpha$.

We train the model of Section~\ref{sec:setup} at each $\alpha$ and then measure, on held-out
sequences and using each sequence's own table $\pi$, the rate at which the model predicts a signal
symbol that is \emph{not} a successor of the current state. We report this separately at
post-distractor positions --- where the state must be recovered by skipping backwards --- and at
post-signal positions, where the state is directly visible.

\begin{table}[!ht]
\caption{Post-distractor error: the rate at which the trained model predicts a signal symbol that
is not a successor of the current state, at positions immediately following a distractor.
Held-out sequences, each scored against its own transition table. The change between
$\alpha=0.7$ and $\alpha=0.8$ is a factor of $143$.}
\label{tab:transition}
\begin{center}
\begin{tabular}{lccccccccc}
\toprule
$\alpha$ & 0.1 & 0.2 & 0.3 & 0.4 & 0.5 & 0.6 & 0.7 & 0.8 & 0.9 \\
\midrule
post-distractor error & 0.89\% & 0.31\% & 0.27\% & 1.16\% & 0.51\% & 0.44\% & 0.36\%
  & \textbf{51.59\%} & \textbf{49.94\%} \\
\bottomrule
\end{tabular}
\end{center}
\end{table}

The transition is abrupt and has no intermediate value in our sweep. Below it the model is
essentially never wrong about which symbols may follow the current state; above it, it is wrong
half the time --- close to what a model that ignores the state entirely and reports the marginal
would achieve.

\paragraph{Interpretation.}
The aggregate cross-entropy does not reveal this. It decreases monotonically in $\alpha$, because
distractors are individually easy to predict (they are uniform on $\mathcal{D}$, and the readout
bias $b_d$ supplies exactly the right mass). The failure is confined to the sub-task that the
aggregate barely weights. A model that has not learned state recovery can still do well overall by
copying the previous symbol, which is the correct state a fraction $\alpha$ of the time; at
$\alpha=0.8$ that shortcut is right four times in five, and the remaining fifth is too small a
share of the loss to force the harder circuit into existence.

% TODO: add the post-signal control column. The claim "state recovery specifically fails" requires
% the post-signal error to stay flat while the post-distractor column explodes. Our current
% measurement of the control had an off-by-one in the state update and is not reportable.
% TODO: seed bands. Table 1 is a single seed.
% TODO: confirm the transition replicates in the reduced-vocabulary configuration used for
% Section 5; Table 1 is measured in the larger configuration (Appendix D).

\section{Learning dynamics from zero initialisation}
\label{sec:dynamics}

To explain Table~\ref{tab:transition} we ask what gradient descent does at the very beginning,
where everything is computable. Initialise $W^{(1)}=W^{(2)}=0$. Both softmaxes are then uniform
over the admissible keys and, by \eqref{eq:arch},
\begin{equation}
h_t=f_t=m_t:=\frac1t\sum_{u<t}e_{v_u},
\end{equation}
the running symbol frequency: the model begins as the in-context unigram predictor.

Split $W^{(2)}$ into $d\times d$ blocks according to $x^{(2)}=[x;h]$, so that the score
decomposes into four interpretable terms,
\begin{equation}
W^{(2)}=\begin{bmatrix}P & Q\\ R & C\end{bmatrix},\qquad
\sqrt d\,S^{(2)}_{t,u}=\underbrace{x_tPx_u^\top}_{\text{symbol--symbol}}
+\underbrace{x_tQh_u^\top}_{\text{induction}}
+\underbrace{h_tRx_u^\top}_{\text{cross}}
+\underbrace{h_tCh_u^\top}_{\text{state--state}}.
\label{eq:blocks}
\end{equation}
The symbol--symbol block $P$ is the in-weights channel: it scores a key by its own symbol against
the query's, which is what a model would use to store a fixed bigram table. The induction block
$Q$ is the in-context channel: it scores a key by the \emph{state} that preceded it.

\begin{theorem}[Dynamics at initialisation]
\label{thm:dynamics}
At $W^{(1)}=W^{(2)}=0$, writing $\omega_t=(m_t[y_t]+b_{y_t})^{-1}>0$:
\begin{enumerate}
\item[(i)] \textup{(Gating.)} $\partial L/\partial W^{(1)}=0$ identically. The first layer cannot
move until the second has left the origin.
\item[(ii)] \textup{(Closed form.)}
$\displaystyle -\frac{\partial L}{\partial W^{(2)}}
=\frac{1}{|\mathcal{T}|\sqrt d}\sum_{t\in\mathcal{T}}\frac{\omega_t}{t}\sum_{u<t}
\bigl(\mathbf 1[v_u=y_t]-m_t[y_t]\bigr)\,x^{(2)}_t\otimes x^{(2)}_u .$
\item[(iii)] \textup{(The in-weights channel is dead.)} In expectation over the resampled table,
the induced update to $P$ has no diagonal component, because
$\mathbb{E}_\pi[\mathbf 1[w=y_t]\mid \mathrm{st}(t)=v]=\alpha/V_s$ independently of
$w\in\mathcal{S}$. The corresponding update to $Q$ does have one.
\end{enumerate}
\end{theorem}

Part (iii) is the reason the title of this section is about in-context learning. The marginal
bigram statistic is uninformative \emph{by construction of the task}: averaged over the resampled
table, every signal symbol is equally likely to follow every other. The model cannot take a step
towards storing a transition table even if such a step would help, because the gradient in that
direction is zero. The first direction available to it is the one that reads the context.

Measured at $\alpha=0.6$, the diagonal-minus-off-diagonal statistic of the first descent
direction, restricted to the signal sub-block, is $-1.48\times10^{-4}$ on $P$ against
$+9.62\times10^{-5}$ on $Q$ and $+9.80\times10^{-6}$ on $C$; the cosine between the induction
block and the construction's is $+0.867$, against a random-matrix baseline of $0.010$. The first
gradient does not merely avoid the wrong direction --- it points most of the way to the right one.

\paragraph{Steps one and two.}
Appendix~\ref{app:dynamics} continues the analysis. After one step the first layer's gradient
becomes non-zero and is proportional to those blocks of the second layer's update that contract
with $h$, so the second layer \emph{recruits} the first
(Lemma~\ref{lem:step1}). After two steps the first layer is driven towards suppressing distractor
keys and sharpening onto the most recent signal position, which is state recovery
(Lemma~\ref{lem:step2}). The circuit is built outwards from the readout rather than inwards from
the input.

\paragraph{Connecting to the transition.}
The coefficient in Theorem~\ref{thm:dynamics}(ii) is supported on positions whose target is a
successor of the current state. The positions at which state recovery is \emph{required} --- those
following a distractor --- form a fraction $1-\alpha$ of the scored window. As $\alpha\to1$ that
share vanishes, and with it the gradient pressure to build the first layer at all.
% TODO: make this quantitative. This is the step that turns Section 5 into an explanation of
% Section 4 rather than a plausible story alongside it.

\section{Reduction to the noiseless case}
\label{sec:reduction}

Our setting contains the standard one as a boundary.

\begin{theorem}[Reduction]
\label{thm:reduction}
As $\alpha\to1$ the distractor sub-vocabulary is never emitted; the state at $t$ is $v_t$ itself;
the construction of Theorem~\ref{thm:construction} degenerates to a previous-token head composed
with an induction head; and the task becomes in-context learning of a Markov chain with a
per-sequence transition table, as studied by \citet{edelman2024}.
\end{theorem}

This positions the prior work as our $\alpha=1$ limit rather than as a competing account, and it
provides a consistency check: our measurements at $\alpha=1.0$ should reproduce the qualitative
behaviour reported there, and they do.
% TODO: state which quantities are compared and give the numbers.

\section{Related work}
\label{sec:related}

\paragraph{Induction heads.}
The induction head was identified as a concrete attention circuit by \citet{elhage2021} and
linked to the emergence of in-context learning in language models by \citet{olsson2022}. Its
canonical form is a two-layer composition: a previous-token head writes the identity of the
preceding symbol into the residual stream, and a second head attends to positions whose written
predecessor matches the current symbol, copying what followed. \citet{bietti2023} study a
two-layer model in which a globally stored component and an in-context component coexist, a
decomposition structurally analogous to our signal/distractor split, though their axis separates
what is memorised from what is inferred, whereas ours separates positions that carry state from
positions that do not.

\paragraph{In-context learning of Markov structure.}
The setting closest to ours is \citet{edelman2024}, who train two-layer attention-only
transformers on sequences whose transition matrix is resampled per sequence and document a
multi-phase trajectory: uniform predictions, then an in-context unigram solution, then a rapid
transition to the in-context bigram solution, with the unigram stage acting as a distractor that
delays the final one. Our model at zero initialisation \emph{is} that unigram solution, which
Theorem~\ref{thm:dynamics} makes precise, and their setting is the $\alpha\to1$ boundary of ours
(Theorem~\ref{thm:reduction}). \citet{ekbote2025} characterise which Markov orders two-layer
transformers can represent, \citet{shi2025} analyse transformers trained on random walks, and
\citet{dai2026} examine whether pre-trained language models infer hidden Markov models in
context; \citet{hao2025} treat the hidden-state case as multi-task feature decoupling. None of
these introduces material that is interleaved with the chain and carries no state.

\paragraph{Training dynamics and the emergence of the circuit.}
\citet{nichani2024} prove that gradient descent recovers the causal structure of a latent graph,
and \citet{unveiling2024} prove a three-stage convergence for a two-layer transformer on $n$-gram
chains in which a feed-forward selector chooses among parents at fixed offsets. \citet{varre2025}
show that sub-$n$-gram solutions are near-stationary points of the in-context $n$-gram objective,
which is the mechanism by which intermediate solutions trap optimisation; \citet{musat2026} and
\citet{wang2025} analyse the emergence and the approximation--dynamics trade-off more broadly, and
\citet{yuksel2026} track the incremental appearance of sparse attention patterns. All
of these analyse architectures in which the informative position lies at a bounded and enumerable
displacement from the query. Our first layer must instead locate a parent at an unbounded,
geometrically distributed lag, and our analysis concerns how that additional primitive is
acquired --- and, for small distractor rates, why it is not.

\paragraph{Selecting causal structure in context.}
\citet{selective2025} study interleaved chains in which the \emph{lag} varies across sequences
while the transition probabilities are held fixed, and construct a three-layer selective induction
head that identifies the operative lag from context. The two settings vary complementary axes:
they vary the dependency distance with the structure fixed, we resample the structure with a
distance that is not selected from a finite menu but must be found by skipping uninformative
material. The mechanisms differ accordingly --- a selector over $k$ candidate lags cannot
represent a geometric lag distribution for any finite $k$, while our first layer expresses it with
a recency ramp and a distractor veto.

\paragraph{Shortcuts, data distribution, and algorithm selection.}
Our empirical finding --- that a capability required on a vanishing fraction of positions is not
acquired, and a shortcut correct on the remainder is acquired instead --- belongs to a line of
work on how the training distribution selects among algorithms a model could in principle
implement. \citet{kawata2025} show that data diversity governs whether a transformer settles on a
shortcut or on a genuine induction head, and \citet{zucchet2025} show how the data distribution
and repetition structure control the emergence of sparse attention. \citet{lavie2026} give a
Bayesian account of phase transitions in copy-head emergence. Relative to these, the axis we vary
is not diversity or repetition but the \emph{density of positions that require a specific
sub-circuit}, and the resulting transition is sharp and non-monotone in task difficulty: the
aggregate objective becomes easier as distractors become rarer, while the sub-task governed by
them becomes catastrophically worse.

\paragraph{Denoising and Bayesian views of attention.}
\citet{smart2025} study in-context denoising with a one-layer transformer, relating attention to
associative memory retrieval; their noise corrupts the observations whereas ours occupies whole
positions and displaces the causal parent. \citet{sun2025} isolate the role of the feed-forward
block in nonlinear in-context learning, which is the component our readout deliberately
replaces with a fixed renormalisation (Appendix~\ref{app:readout}). Bayesian treatments of attention
\citep{agarwal2026,wakayama2025}, of latent mixture models learned by mirror descent
\citep{dangelo2026}, and of selective state-space models \citep{zhang2026} supply the normative
reference against which we measure: our construction
realises the context-Bayes predictor under add-$\beta$ smoothing, and the phase transition is a
failure of optimisation to reach it, not a failure of the architecture to express it.

\paragraph{Parameterisation and symmetry.}
Weight-space comparisons between a learned circuit and a hand construction are only meaningful
modulo the architecture's exact symmetries. \citet{silverstein2026} quantify how much of an
optimiser's motion is spent on such gauge directions rather than on function change, and
\citet{kobayashi2024} show that weight decay drives attention layers towards low rank, which is a
further implicit restriction of the reachable set. Appendix~\ref{app:gauge} records the symmetries
of our architecture and uses them to separate the freezing choices that cost generality from those
that do not.

\section{Discussion and limitations}
\label{sec:limitations}

The phenomenon we report is a failure of a rare sub-task to attract gradient, and the mechanism is
not specific to distractors: whenever a capability is required on a vanishing fraction of the
training distribution, and a shortcut exists that is correct on the remainder, the shortcut is
what training will produce. What makes the present setting useful is that the rare sub-task, the
shortcut and the correct circuit are all explicitly identifiable.

Several limitations should be stated plainly. The task is synthetic, the vocabulary small and the
model two layers deep; we make no claim about trained language models. The value circuits are
held fixed, which Appendix~\ref{app:gauge} shows is free in the first layer and a real restriction
in the second. Theorem~\ref{thm:dynamics}(iii) is proved for the expectation of the in-weights
block; the corresponding lower bound for the in-context block is verified numerically but its
proof requires control of the chain's mixing, which we do not yet have. Finally, weight-space
agreement with the construction is not monotone during training: it peaks within the first ten
steps and then declines while the loss continues to improve. We report this but cannot explain it.


\subsubsection*{Reproducibility statement}

The data-generating process is specified completely in Section~\ref{sec:setup}; the construction
of Theorem~\ref{thm:construction} is given in closed form in Appendix~\ref{app:construction} and
requires no fitting. Architecture, optimiser, initialisation, the two configurations used and the
evaluation protocol for Table~\ref{tab:transition} are in Appendix~\ref{app:exp}. Proofs of
Theorem~\ref{thm:dynamics} are in Appendix~\ref{app:dynamics}, with the one step that is currently
verified numerically rather than proved marked explicitly there. Source code, the data-generation
script and the checkpoints used for every reported number are included in the supplementary
material.
% TODO: confirm the supplementary actually contains these before submitting.

\subsubsection*{AI use statement}

% MANDATORY for ICLR 2027; does not count toward the page limit.
% TODO -- THE AUTHORS MUST EDIT THIS TO MATCH WHAT ACTUALLY HAPPENED AND TAKE RESPONSIBILITY
% FOR IT. The draft below reflects our understanding of how the work was carried out; please
% correct anything that is inaccurate. ICLR treats a substantial falsehood produced by an LLM as
% a Code of Ethics violation, so every number, citation and claim in the final text must be
% checked by the authors by hand.
We used an AI assistant for the following tasks: implementing the training and evaluation code;
executing numerical experiments and tabulating their outputs; searching for and summarising
related literature; and drafting portions of this manuscript. The research question, the task
design, the construction of Section~\ref{sec:construction} and the theoretical claims were
developed by the authors. All numerical results reported here were reproduced by the authors from
the released code, and all cited works were consulted by the authors in the original. The authors
take full responsibility for the content of this paper, including any errors.

\bibliography{iclr2027_conference}
\bibliographystyle{iclr2027_conference}

\appendix

\section{Learning dynamics from zero initialisation: proofs}
\label{app:dynamics}

\subsection{Notation}
\label{app:notation}

Let $\mathcal{S}=\{0,\dots,V_s-1\}$ be the signal symbols and
$\mathcal{D}=\{V_s,\dots,V-1\}$ the distractors, $V=V_s+V_d$. A sequence is
$v_{0:T-1}$ with $y_t=v_{t+1}$ the target at position $t$; the loss is averaged over the scored
window $t\in\mathcal{T}=[T/2,T)$. For each sequence a transition table $\pi$ is drawn
independently, assigning to every $s\in\mathcal{S}$ a successor set
$\mathrm{succ}_\pi(s)\subset\mathcal{S}$ of size two. Writing $\mathrm{st}(t)$ for the chain state
after position $t$, the generative rule is: with probability $\alpha$ emit a uniform element of
$\mathrm{succ}_\pi(\mathrm{st}(t-1))$ and advance the state; with probability $1-\alpha$ emit a
uniform element of $\mathcal{D}$ and leave the state unchanged.

The input embedding is the concatenation of two one-hots,
$x_t=[\,e_{v_t}\,;\,e_t\,]\in\mathbb{R}^{d}$ with $d=V+T$; there is no learned embedding. Both
layers use a strictly causal mask $M_{t,u}=-\infty$ for $u\ge t$. With
$W^{(1)}\in\mathbb{R}^{d\times d}$ and $W^{(2)}\in\mathbb{R}^{2d\times 2d}$ the combined
query--key matrices,
\begin{align}
S^{(1)}_{t,u} &= \tfrac{1}{\sqrt d}\,x_t W^{(1)} x_u^\top + M_{t,u},
& a^{(1)}_{t,\cdot} &= \mathrm{softmax}\bigl(S^{(1)}_{t,\cdot}\bigr),
& h_t &= \sum_{u} a^{(1)}_{t,u}\,e_{v_u}, \\
x^{(2)}_t &= [\,x_t\,;\,h_t\,]\in\mathbb{R}^{2d},
& S^{(2)}_{t,u} &= \tfrac{1}{\sqrt d}\,x^{(2)}_t W^{(2)} x^{(2)\top}_u + M_{t,u},
& f_t &= \sum_{u} a^{(2)}_{t,u}\,e_{v_u},
\end{align}
where we have used that the value circuits are fixed to the selector that passes the symbol
one-hot (Appendix~\ref{app:gauge}). Note that $f_t$ depends on $h$ \emph{only} through the
scores $S^{(2)}$. The readout is a renormalisation with a fixed bias $b$,
\begin{equation}
p_t=\frac{f_t+b}{\textstyle\sum_w (f_t[w]+b_w)},\qquad
L=-\frac{1}{|\mathcal{T}|}\sum_{t\in\mathcal{T}}\log p_t[y_t].
\end{equation}
Because $f_t$ is a convex combination of one-hots, $\sum_w f_t[w]=1$ identically, so the
denominator $S:=1+V_s b_s+V_d b_d$ is a constant independent of the context.

It is convenient to split $W^{(2)}$ into $d\times d$ blocks according to
$x^{(2)}=[x;h]$:
\begin{equation}
W^{(2)}=\begin{bmatrix} P & Q \\ R & C\end{bmatrix},
\qquad
S^{(2)}_{t,u}\sqrt d = \underbrace{x_tPx_u^\top}_{\text{symbol--symbol}}
+\underbrace{x_tQh_u^\top}_{\text{induction}}
+\underbrace{h_tRx_u^\top}_{\text{cross}}
+\underbrace{h_tCh_u^\top}_{\text{state--state}} .
\end{equation}

\subsection{Lemma 1: the gradient at initialisation}
\label{app:lemma1}

\begin{lemma}[Step $0$]
\label{lem:step0}
Let $W^{(1)}=W^{(2)}=0$. Then both attention distributions are uniform on the admissible keys,
$a^{(1)}_{t,u}=a^{(2)}_{t,u}=1/t$ for $u<t$, and consequently
\begin{equation}
h_t=f_t=m_t:=\frac{1}{t}\sum_{u<t}e_{v_u},
\end{equation}
the running symbol frequency of the context. Writing
$\omega_t:=\bigl(m_t[y_t]+b_{y_t}\bigr)^{-1}>0$, the gradients are
\begin{align}
-\frac{\partial L}{\partial W^{(2)}}
&=\frac{1}{|\mathcal{T}|\sqrt d}\sum_{t\in\mathcal{T}}\frac{\omega_t}{t}
\sum_{u<t}\Bigl(\mathbf 1[v_u=y_t]-m_t[y_t]\Bigr)\;
x^{(2)}_t\otimes x^{(2)}_u ,
\label{eq:g2}\\[2pt]
\frac{\partial L}{\partial W^{(1)}}&=0 .
\label{eq:g1}
\end{align}
\end{lemma}

\begin{proof}
At $W^{(1)}=W^{(2)}=0$ every score is $0$, so each softmax is uniform over $\{u<t\}$ and the
first display follows. For \eqref{eq:g2}, treat $f_t$ as a free variable and set
$g_t[w]:=\partial L_t/\partial f_t[w] = -\delta_{w,y_t}\omega_t + 1/S$. Since
$f_t=\sum_u a^{(2)}_{t,u}e_{v_u}$, the softmax Jacobian gives
\begin{equation}
\frac{\partial L_t}{\partial S^{(2)}_{t,u}}
= a^{(2)}_{t,u}\Bigl(g_t[v_u]-\sum_{u'}a^{(2)}_{t,u'}g_t[v_{u'}]\Bigr)
= a^{(2)}_{t,u}\bigl(g_t[v_u]-g_t\!\cdot\! f_t\bigr).
\end{equation}
Using $\sum_w f_t[w]=1$ we get $g_t\!\cdot\!f_t=-m_t[y_t]\omega_t+1/S$, whence
$g_t[v_u]-g_t\!\cdot\!f_t=-\omega_t(\mathbf 1[v_u=y_t]-m_t[y_t])$. Composing with
$\partial S^{(2)}_{t,u}/\partial W^{(2)}=x^{(2)}_t\otimes x^{(2)}_u/\sqrt d$ gives \eqref{eq:g2}.

For \eqref{eq:g1}, observe that $h$ enters the loss only through $S^{(2)}$, and
\begin{equation}
\frac{\partial S^{(2)}_{t,u}}{\partial h_u}=\frac{1}{\sqrt d}\bigl(x^{(2)}_tW^{(2)}\bigr)_{d:2d},
\qquad
\frac{\partial S^{(2)}_{t,u}}{\partial h_t}=\frac{1}{\sqrt d}\bigl(W^{(2)}x^{(2)\top}_u\bigr)_{d:2d},
\end{equation}
both of which carry $W^{(2)}$ as a multiplicative factor and therefore vanish at
$W^{(2)}=0$. Since every path from $W^{(1)}$ to $L$ passes through $h$, the chain rule gives
$\partial L/\partial W^{(1)}=0$.
\end{proof}

Equation \eqref{eq:g2} has a direct reading: the first step raises the score of exactly those
keys whose \emph{symbol equals the target}, relative to a baseline $m_t[y_t]$. Equation
\eqref{eq:g1} is a gating statement --- the two layers cannot form simultaneously.

\begin{corollary}[The in-weights channel receives no signal]
\label{cor:bigram}
Decompose \eqref{eq:g2} into the blocks of Section~\ref{app:notation}. The symbol--symbol block
is
\begin{equation}
P^\star[v,w]\;\propto\;\sum_{t:\,v_t=v}\frac{\omega_t}{t}
\sum_{u<t:\,v_u=w}\bigl(\mathbf 1[w=y_t]-m_t[y_t]\bigr).
\end{equation}
Taking expectations over the resampled table $\pi$, for $w\in\mathcal{S}$,
\begin{equation}
\mathbb{E}_\pi\bigl[\mathbf 1[w=y_t]\,\big|\,\mathrm{st}(t)=v\bigr]
=\frac{\alpha}{2}\,\mathbb{P}_\pi\bigl[w\in\mathrm{succ}_\pi(v)\bigr]
=\frac{\alpha}{V_s},
\end{equation}
which is independent of $w$. Hence $\mathbb{E}[P^\star]$ has no diagonal component: the marginal
bigram statistic carries no information, and the in-weights solution is not merely suboptimal but
\emph{gradient-inaccessible} at initialisation.
\end{corollary}

By contrast the induction block $Q^\star[v,\cdot]$ pairs the coefficient
$\mathbf 1[v_u=y_t]-m_t[y_t]$ with the key feature $h_u=m_u$. Conditioned on $v_u=w$ with
$w\in\mathrm{succ}_\pi(v)$, occurrences of $w$ are generated by transitions out of the
predecessors of $w$, of which $v$ is one; therefore $m_u[v]$ is elevated relative to its
unconditional value, and the elevation occurs exactly on the diagonal $v=v$. This is the
in-context correlation that survives the average over $\pi$.

\paragraph{Status of Corollary~\ref{cor:bigram}.}
% TODO -- THE MAIN REMAINING TECHNICAL WORK.
The first display of the corollary is exact. The concluding paragraph is a mechanism sketch:
turning ``$m_u[v]$ is elevated'' into a quantitative lower bound on the diagonal of
$\mathbb{E}[Q^\star]$ requires control of the chain's mixing, since $m_u$ is an average over the
whole prefix rather than over the single position preceding $u$. We currently verify the
conclusion numerically. At $\alpha=0.6$, single seed, the diagonal-minus-off-diagonal statistic
of the first descent direction restricted to the signal sub-block is
$-1.48\times10^{-4}$ for $P^\star$, $+9.62\times10^{-5}$ for $Q^\star$ and
$+9.80\times10^{-6}$ for $C^\star$; the cosine against the construction is $+0.867$ on the
induction block against a random-matrix baseline of $0.010$.

\subsection{Lemma 2: the first layer unlocks}
\label{app:lemma2}

\begin{lemma}[After step $1$]
\label{lem:step1}
Let $W^{(2)}_1=\eta G^{(2)}$ be the result of one gradient step from Lemma~\ref{lem:step0}, with
$G^{(2)}=-\partial L/\partial W^{(2)}$, and $W^{(1)}_1=0$ by \eqref{eq:g1}. Then:
\begin{enumerate}
\item[(i)] the first layer is unchanged, so $h_t=m_t$ still;
\item[(ii)] the second layer's attention is perturbed to first order,
$a^{(2)}_{t,u}=\tfrac1t\bigl(1+\tfrac{\eta}{\sqrt d}(\sigma_{t,u}-\bar\sigma_t)\bigr)+O(\eta^2)$
with $\sigma_{t,u}=x^{(2)}_tG^{(2)}x^{(2)\top}_u$ and $\bar\sigma_t=\tfrac1t\sum_{u<t}\sigma_{t,u}$;
\item[(iii)] the first layer's gradient becomes non-zero and is of order $\eta$:
\begin{equation}
-\frac{\partial L}{\partial W^{(1)}}\Big|_{1}
=\eta\,\Phi\bigl(G^{(2)}\bigr)+O(\eta^{2}),
\end{equation}
where $\Phi$ is linear in the blocks of $G^{(2)}$ that contract with $h$, namely $Q,R$ and $C$.
\end{enumerate}
\end{lemma}

\begin{proof}
(i) is \eqref{eq:g1}. (ii) is a first-order expansion of the softmax about the uniform
distribution. For (iii), the derivatives of $S^{(2)}$ with respect to $h$ displayed in the proof
of Lemma~\ref{lem:step0} are now evaluated at $W^{(2)}=\eta G^{(2)}$ rather than $0$, so they are
$\eta$ times the corresponding blocks of $G^{(2)}$; composing with
$\partial h/\partial W^{(1)}$, which is $O(1)$ and non-zero, gives the claim. Only the blocks
that multiply $h$ survive the contraction, which are $Q$, $R$ and $C$.
\end{proof}

The content of Lemma~\ref{lem:step1} is that the second layer \emph{recruits} the first. The
first layer receives gradient only in proportion to how much the second layer has already
learned to read $h$, and the direction it receives is dictated by the blocks of $G^{(2)}$ that
consume $h$ --- predominantly $Q^\star$, whose diagonal structure was established above.

% TODO: the Adam caveat. Our experiments use Adam, whose first update has magnitude set by the
% learning rate and is essentially independent of the gradient scale. The O(eta) bookkeeping
% above is therefore the gradient-flow / plain-GD statement. State which regime the empirical
% section reports, or repeat the measurement under plain GD.

\subsection{Lemma 3: what the first layer is pushed to compute}
\label{app:lemma3}

\begin{lemma}[After step $2$]
\label{lem:step2}
To leading order the update to $W^{(1)}$ increases $S^{(1)}_{u,s}$ for pairs $(u,s)$ with $s<u$
such that $v_s\in\mathcal{S}$, and decreases it for $v_s\in\mathcal{D}$; among admissible $s$ the
increase is monotone in $s$. Consequently $h_u$ moves from the uniform prefix average $m_u$
towards the indicator of the most recent signal symbol before $u$.
\end{lemma}

% TODO: PROOF NOT YET WRITTEN. The intended argument:
%   * By Lemma 2(iii) the layer-1 gradient is Phi(G^{(2)}), dominated by the Q block.
%   * Q^star[v, .] rewards keys u whose h_u places mass on v, when v is the query symbol.
%   * h_u places mass on v in proportion to how often v occurs in the prefix of u; the update
%     therefore sharpens h_u onto whichever symbol best predicts v_u, which by the generative
%     process is the state, i.e. the most recent signal symbol.
%   * The distractor suppression and the recency ordering come from the two terms of the
%     coefficient (indicator minus baseline).
% This is the step that completes the story and it is the second piece of remaining work.

\paragraph{Summary of the three steps.}
At step $0$ the model is the in-context unigram predictor; the first layer is frozen by
\eqref{eq:g1}, and the only structured gradient available to the second layer lies in the
induction block, because the resampled table annihilates the marginal bigram
(Corollary~\ref{cor:bigram}). At step $1$ the second layer has moved, which unlocks the first
(Lemma~\ref{lem:step1}). At step $2$ the first layer is driven towards state recovery
(Lemma~\ref{lem:step2}). The circuit is therefore built outwards from the readout, not inwards
from the input, and the first direction it can take is already an in-context one.

% TODO: connect to Section 4. The coefficient in \eqref{eq:g2} is supported on positions whose
% target is a successor; the share of the scored window on which state recovery is *required*
% -- the post-distractor positions -- is proportional to (1 - alpha). Making this quantitative
% is what converts the dynamics into an explanation of the phase transition.

\section{Proof of Theorem~\ref{thm:construction}}
\label{app:construction}

We exhibit weights and verify the two claims. Write $\lambda>0$ for a recency slope, $K>0$ for a
match bonus and $B>0$ for a veto depth; all three are free and only their ratios matter.

\paragraph{Layer 1 (state recovery).}
Set $W^{(1)}$ so that the score decomposes as
\begin{equation}
\sqrt d\,S^{(1)}_{t,u}\;=\;\underbrace{-B\,\mathbf 1[v_u\in\mathcal{D}]}_{\text{distractor veto}}
\;\underbrace{-\,\lambda\,(t-u)}_{\text{recency}} .
\label{eq:l1score}
\end{equation}
Concretely the symbol rows of the key side carry $-B$ on the distractor coordinates and the
position block carries $-\lambda(t-u)$; the query side is a scaled identity, so the
$\sqrt d$ cancels. Among admissible keys ($u<t$) the score is maximised by the largest $u$ with
$v_u\in\mathcal{S}$. Two candidates separated by one position differ by $\lambda$, so for
$B\gg\lambda\gg1$ the softmax places mass $1-O(e^{-\lambda})$ on the most recent signal position
and $O(e^{-B})$ on every distractor. Hence $h_t=e_{\mathrm{st}(t)}$ up to those errors, and
crucially the lag is unbounded: \eqref{eq:l1score} is linear in $t-u$ and so ranks correctly at
every separation with a fixed number of parameters.\footnote{A capped ramp,
$-\lambda\min(t-u,c)$, behaves identically because only the argmax matters, and is what our
trained models approximate.}

\paragraph{Layer 2 (counting induction).}
Set $W^{(2)}$ so that
\begin{equation}
\sqrt d\,S^{(2)}_{t,u}\;=\;K\,\mathbf 1\bigl[h_u=e_{v_t}\bigr]
\;+\;\tfrac{K}{20}\,\mathbf 1\bigl[h_u=h_t\bigr]\;-\;B\,\mathbf 1[v_u\in\mathcal{D}] .
\label{eq:l2score}
\end{equation}
The first term is the induction match: it rewards keys whose \emph{recovered state} equals the
query's own symbol. The second is the same question asked of the query's recovered state, and
takes over when $v_t$ is a distractor, in which case the first term is identically zero because
$e_{v_t}$ has no support on $\mathcal{S}$. The third removes distractor keys from the average.

Let $A_t=\{u<t: v_u\in\mathcal{S},\ \mathrm{st}(u-1)=\mathrm{st}(t)\}$ be the set of past
positions whose predecessor state equals the current state. By \eqref{eq:l2score} every
$u\in A_t$ scores $K(1+\tfrac1{20})$ and every other admissible $u$ scores at most $0$, so
$a^{(2)}_{t,u}=\mathbf 1[u\in A_t]/|A_t|+O(e^{-K})$ and
\begin{equation}
f_t[w]=\sum_u a^{(2)}_{t,u}\mathbf 1[v_u=w]
=\frac{\#\{u\in A_t: v_u=w\}}{|A_t|}+O(e^{-K}),
\end{equation}
which is the claimed empirical successor frequency. Two consequences are worth recording:
$\sum_w f_t[w]=1$ exactly, and $f_t[w]=0$ for $w\in\mathcal{D}$, since $A_t$ contains only signal
positions. Composing with \eqref{eq:readout} and Appendix~\ref{app:readout} gives add-$b_s$
smoothed context-Bayes prediction. \qed

\section{The readout, and why no softmax is needed}
\label{app:readout}

\paragraph{Derivation of $b_d$.}
Write $D=\sum_w(f_t[w]+b_w)$. A model with a perfect distractor filter has $f_t$ supported on
$\mathcal{S}$ with $\sum_w f_t[w]=1$, so $D=1+V_sb_s+V_db_d$, a constant independent of the
context.\footnote{This is the property that makes a single scalar bias sufficient. It fails if the
value circuit is trained, since then $\sum_w f_t[w]$ varies with the context and no fixed $b_d$
delivers a fixed distractor mass.} Demanding that each distractor receives exactly
$(1-\alpha)/V_d$,
\begin{equation}
\frac{b_d}{D}=\frac{1-\alpha}{V_d}
\;\Longrightarrow\;
V_db_d=(1-\alpha)D
\;\Longrightarrow\;
D=\bigl(1+V_sb_s\bigr)+(1-\alpha)D
\;\Longrightarrow\;
D=\frac{1+V_sb_s}{\alpha},
\end{equation}
whence $b_d=(1-\alpha)(1+V_sb_s)/(V_d\alpha)$ as in \eqref{eq:readout}, and the signal side
receives $p_t[w]=\alpha(f_t[w]+b_s)/(1+V_sb_s)$, exactly proportional to the smoothed count.

\paragraph{Relation to a log-plus-softmax readout.}
A conventional readout would apply an unembedding and a softmax. If the intervening nonlinearity
is $z_w=\log(f_t[w]+b_s)$ on signal coordinates and a constant $\beta$ on distractor coordinates,
then $\mathrm{softmax}(z)_w=(f_t[w]+b_s)/\sum_{w'}e^{z_{w'}}$, which is \eqref{eq:readout} with
$b_d=e^{\beta}$. The logarithm and the exponential cancel and the pair implements precisely the
renormalisation we use.\footnote{Measured difference between the two implementations on our
checkpoints: $5.6\times10^{-17}$.} The logarithm is not optional decoration: without it the
softmax would return $e^{f_t[w]}$, exponential rather than linear in the count, and no rescaling
repairs this because $\mathrm{softmax}(\kappa f)$ is a different function of $f$ than
$f/\sum f$ for every $\kappa$.

\section{Experimental details}
\label{app:exp}

\paragraph{Data.}
Sequences are generated as in Section~\ref{sec:setup} with successor sets drawn uniformly without
replacement and with equiprobable successors. Train and test tables are drawn independently; no
table is shared between splits, so test performance measures in-context inference and not table
memorisation.

\paragraph{Two configurations.}
% TODO: this is a real scope issue and should be resolved before submission.
Table~\ref{tab:transition} is measured in the larger configuration
($V_s=V_d=10$, $T=200$, scored on $[100,200)$) with a trained unembedding and the log-plus-softmax
readout of Appendix~\ref{app:readout}. The dynamics of Section~\ref{sec:dynamics} are measured in
a reduced configuration ($V_s=V_d=4$, $T=100$, scored on $[50,100)$) with the fixed readout
\eqref{eq:readout}, which is what makes the closed forms of Theorem~\ref{thm:dynamics} available.
The two configurations should be reconciled: either the phase transition should be replicated in
the reduced configuration, or the dynamics should be re-derived for the larger one.

\paragraph{Optimisation.}
Full-batch Adam, constant learning rate, zero initialisation of both query--key matrices, with the
value circuits and readout fixed throughout. Checkpoint selection is on held-out loss.

\paragraph{Reported numbers.}
Table~\ref{tab:transition} uses 1500 held-out sequences per $\alpha$ and scores every position in
the window whose state is defined. The gradient statistics in Section~\ref{sec:dynamics} are exact
full-batch gradients at initialisation, not stochastic estimates.
% TODO: all reported numbers are a single seed. Add seed bands before submission.

\section{Weight-space agreement is not monotone}
\label{app:peak}

We record an observation we cannot yet explain. Measuring the agreement between the trained
$W^{(2)}$ and the construction of Theorem~\ref{thm:construction}, after quotienting the exact
symmetries of Appendix~\ref{app:gauge}, the agreement rises to its maximum within the first ten
steps and then \emph{declines}, while the loss continues to improve:

\begin{center}
\begin{tabular}{lcccc}
\toprule
& peak agreement & at step & final agreement & loss over that interval \\
\midrule
$\alpha=0.6$ & 0.911 & 5 & 0.779 & $1.950\to1.698$ \\
$\alpha=0.8$ & 0.885 & 7 & 0.740 & $1.690\to1.384$ \\
\bottomrule
\end{tabular}
\end{center}

Gradient descent reaches the construction almost immediately and then trades roughly $0.14$ of
agreement for roughly $0.3$ nats. The construction is therefore a way-station and not the
destination, and characterising what is gained in that trade is open.

\section{Why the value circuit may be held fixed}
\label{app:gauge}

% TODO. The asymmetry, which justifies the setup in Section 5:
%  * For any invertible A, transforming the first layer's value matrix by A and compensating in
%    the second layer's query-key and value matrices leaves every output bit-identical.
%    Fixing the first layer's value circuit is therefore a choice of coordinates, not an
%    assumption. (Verified: cross-entropy unchanged to 0.00e+00 under a random A.)
%  * No perturbation of the second layer's value circuit admits a compensating change in the
%    query-key matrices, because the former sits outside the attention softmax and the latter
%    inside it. Fixing it is a genuine restriction, and we state it as such.

\end{document}
